\documentclass[11pt]{article}
\usepackage[a4paper,margin=25mm]{geometry}
\usepackage{amsmath,amssymb,amsthm,mathtools,microtype}
\usepackage[T1]{fontenc}\usepackage{lmodern}
\usepackage[colorlinks=true,linkcolor=blue,urlcolor=blue]{hyperref}
\newtheorem{theorem}{Theorem}[section]
\newtheorem{lemma}[theorem]{Lemma}
\newtheorem{corollary}[theorem]{Corollary}
\newcommand{\F}{\mathbb F}
\newcommand{\Span}{\operatorname{span}}
\newcommand{\SV}{\operatorname{SV}}
\newcommand{\GL}{\operatorname{GL}}
\title{Constant Row Spaces and Low-Degree Concentration\\
for Invertible Affine ROABP Words}
\author{icekylinx \and GPT-6 Astra Pro}\date{26 October 2026}
\begin{document}\maketitle
\begin{abstract}
For an invertible-local width-\(w\) polynomial matrix word, translating
each latent variable by any element transcendental over the base field
makes its row coefficient space generated in total degree at most
\(w-1\). The shifts may be correlated.
The proof tracks constant coefficient row spaces and charges local
derivatives to dimension increases. It works in characteristic zero
or in characteristic greater than the local entry-degree bound.
Consequently \(\gamma+\SV_{w-1}\), with \(\gamma_i(t)=t^i\),
hits unknown invertible affine orbits of words whose local ranks
are one or \(w\). The map has \(2w-1\) seeds; pure \(\SV_w\)
also hits. The argument does not supply an all-characteristic
affine transport theorem.
\end{abstract}

\noindent\textbf{Status.}
This proof has independent model-assisted internal reconstructions
and an exact-source audit. This is not external peer review,
and no publication-priority claim is made.

\section{Statement}

Let \(n,w,\delta\ge1\), and let
\(A_i(T)\in\F[T]^{w\times w}\) have entry degrees at most
\(\delta\), with constant scalar endpoints \(u,v\).
Assume \(\operatorname{char}\F=0\) or
\(\operatorname{char}\F=p>\delta\).
Fix \(L\in\GL_n(\F)\) with coordinate forms \(\ell_i\), and
constants \(c_i\). The representation and affine change are unknown.
Consider
\[
 f(X)=uA_1(\ell_1(X)+c_1)\cdots A_n(\ell_n(X)+c_n)v.
\]

Choose \(n\) distinct abscissas and their Lagrange polynomials
\(\lambda_i(y)\). Let
\[
 \gamma_i(t)=t^i,\qquad
 \mathcal H_i(t,\boldsymbol y,\boldsymbol z)
  =t^i+\sum_{j=1}^{w-1}z_j\lambda_i(y_j).
\]
For \(w=1\), the sum is empty.

\begin{theorem}\label{thm:hit}
If each \(A_i\) has generic rank one or \(w\), then
\[
 f\ne0\quad\Longrightarrow\quad f(\mathcal H)\ne0.
\]
The standard map \(\SV_w\) also hits this class.
With \(D=n\delta\), a rectangular grid for \(\mathcal H\)
uses at most
\[
 (nD+1)\bigl((n-1)D+1\bigr)^{w-1}(D+1)^{w-1}
\]
queries, over a working field of at least \(nD+1\) elements.
Appropriate extension-field oracle evaluations are required if
the base field is too small.
\end{theorem}

The claim is polynomial for fixed \(w\).
It includes arbitrary invertible-local words. It does not assert the
same map for intermediate singular ranks.
Here \(\delta\) bounds the local matrix entries in a representation
with the stated rank pattern; truncating a representation from a
scalar-degree promise alone need not preserve that pattern.
When \(n=0\) or \(\delta=0\), the scalar is constant and one evaluation
suffices. The width-one case in fact works in every characteristic,
since it is a product of nonzero univariate factors at transcendental
arguments.

\section{A univariate block-jet lemma}

For a polynomial row matrix \(R(T)\), its \emph{constant coefficient
row space} is the \(\F\)-span of every row of every coefficient
matrix of \(R\). It is a subspace of \(\F^w\).

\begin{lemma}\label{lem:block}
Let \(B\in\F^{r\times w}\) have row rank \(r\), and let
\(A(T)\in\F[T]^{w\times w}\) be generically invertible with
entry degrees at most \(\delta\).
Let \(U\subseteq\F^w\) be the constant coefficient row space
of \(BA(T)\), and put \(s=\dim U\).
Then \(s\ge r\), and over \(\F(T)\) the rows of
\[
 BA^{(a)}(T),\qquad 0\le a\le s-r,
\]
span \(U\otimes_\F\F(T)\).
The same statement holds after substituting any \(C\) transcendental
over \(\F\), in any extension field containing \(C\).
\end{lemma}
\begin{proof}
Let
\[
 V_k=\Span_{\F(T)}
 \{\text{rows of }BA^{(a)}(T):0\le a\le k\}.
\]
Since \(A(T)\) is invertible, \(\dim V_0=r\).
All these rows lie in \(U\otimes\F(T)\).

If \(V_k=V_{k+1}\), differentiating every displayed spanning row
of \(V_k\) stays in \(V_{k+1}=V_k\).
Differentiating a rational linear combination also stays in \(V_k\),
by the ordinary product rule. Thus \(V_k\) is derivative-stable
and contains all higher derivatives of \(BA\).

These derivatives span \(U\otimes\F(T)\).
Indeed \(\delta!\) is nonzero under the characteristic hypothesis,
and the exact Taylor identity
\[
 BA(S)=\sum_{a=0}^{\delta}
       \frac{BA^{(a)}(T)}{a!}(S-T)^a
\]
holds with an independent variable \(S\).
It expresses every constant coefficient row of \(BA(S)\)
in their \(\F(T)\)-span. The opposite inclusion is immediate.

Therefore the filtration cannot stall before it reaches dimension
\(s\). Starting at dimension \(r\), it is full by \(k=s-r\).
In particular \(s\ge r\).

Choose an \(s\)-by-\(s\) nonzero minor certifying this spanning rank.
Its entries are univariate polynomials over \(\F\), so the minor
remains nonzero at a transcendental \(C\).
The evaluated rows remain contained in \(U\) after scalar extension,
and therefore span it. Further coefficient-field extension does not
change this rank.
\end{proof}

Only derivatives through the local degree \(\delta\) are needed.
Thus positive characteristic requires \(p>\delta\), not
\(p>w\) or \(p>n\delta\). Vanishing higher derivatives cause
no difficulty in the proof.

\section{Constant prefix spaces}

Let \(u\in\F^{1\times w}\) be nonzero and assume all local
matrices are generically invertible.
For each \(j\), define
\[
 R_j(T_1,\ldots,T_j)=uA_1(T_1)\cdots A_j(T_j),
\]
and let \(U_j\subseteq\F^w\) be its constant coefficient row space.
Put \(r_j=\dim U_j\), so \(U_0=\Span_\F(u)\) and \(r_0=1\).
These spaces use \emph{independent original variables}.
They are not coefficient spans of a correlated curve specialization.

If \(B\) is any constant basis matrix of \(U_{j-1}\), the coefficient
row space of \(BA_j(T)\) is exactly \(U_j\). This follows by extracting
coefficients in the independent first \(j-1\) variables.
Consequently
\[
 1=r_0\le r_1\le\cdots\le r_n\le w
\]
by Lemma~\ref{lem:block}.

\begin{theorem}[Low total-degree row concentration]\label{thm:conc}
Let \(K\supseteq\F\) be any extension field, and let
\(C_1,\ldots,C_n\in K\) each be transcendental over \(\F\).
They need not be algebraically independent.
For every \(j\), the coefficient rows of
\[
 R_j(C_1+Y_1,\ldots,C_j+Y_j)
\]
of total degree at most \(r_j-1\) span \(U_j\otimes_\F K\).
\end{theorem}
\begin{proof}
The assertion for \(j=0\) is immediate.
Suppose it is true at \(j-1\). Write \(r=r_{j-1}\) and \(s=r_j\).
The selected prefix coefficient rows of degree at most \(r-1\)
span every row of a constant basis \(B\) of \(U_{j-1}\), after
extending scalars to \(K\).

Lemma~\ref{lem:block}, evaluated at \(C_j\), says that
\[
 BA_j^{(a)}(C_j),\qquad 0\le a\le s-r,
\]
span \(U_j\otimes K\).
When \(a\le\delta\), division by \(a!\) is allowed, and
these are scalar multiples of the corresponding Hasse coefficients
of \(BA_j(C_j+Y_j)\). Terms with \(a>\delta\) are zero and may
be ignored.

Multiply the previous low-degree coefficient rows by these local
Hasse matrices. The products are actual coefficient rows of the
new prefix, with total degree at most
\[
 (r-1)+(s-r)=s-1.
\]
Their \(K\)-span contains the displayed full spanning family,
and hence is \(U_j\otimes K\).
\end{proof}

The rank certificate used at step \(j\) is a polynomial in the
single local argument over \(\F\), not a polynomial involving the
earlier \(C_i\). This is why correlations among the shifts do
not invalidate the induction.

\begin{corollary}\label{cor:scalar}
If the invertible-local scalar \(uA_1(T_1)\cdots A_n(T_n)v\)
is nonzero, then after the shift \(T_i=C_i+Y_i\) its first
nonzero homogeneous term has ordinary total degree at most \(w-1\).
\end{corollary}
\begin{proof}
Nonzeroness says that \(v\) does not annihilate \(U_n\).
The spanning coefficients in Theorem~\ref{thm:conc} therefore
include one whose pairing with \(v\) is nonzero.
\end{proof}

In characteristic zero the same concentration statement holds for
generically invertible rational univariate local matrices.
Define their constant row spaces by clearing scalar denominators.
The choice of clearing factor does not affect the row space:
a constant column annihilates a rational row if and only if it
annihilates any nonzero scalar multiple of that row.
Clear a nonzero scalar univariate denominator from each local
matrix. The polynomial argument applies with its finite local degree.
At every transcendental \(C_i\), the cleared scalar is a unit in
the residual power-series ring. Multiplication by a scalar unit
relates coefficient rows up to any total degree by an invertible
triangular transformation, preserving the asserted concentration.
This rational extension does not assert the same \(p>\delta\)
bound when clearing denominators raises local degrees.

\section{Affine detection and rank-one factors}

For the curve \(\gamma_i=t^i\), every
\[
 C_i(t)=c_i+\ell_i(\gamma(t))
\]
is a nonconstant polynomial in \(t\), hence transcendental over \(\F\).
This holds in every characteristic. Ordinary derivatives of the
\(C_i\) need not be nonzero; the proof uses transcendence only.

For an invertible-local word, Corollary~\ref{cor:scalar} gives a
nonzero first homogeneous term of degree at most \(w-1\).
The invertible residual change \(Y=LX\) preserves that degree
and nonzeroness.
Every nonzero polynomial of total degree at most \(w-1\) is hit by
\(\SV_{w-1}\): some nonzero monomial uses at most \(w-1\)
coordinates, and this map contains every such coordinate plane
by specializing its direction seeds to the corresponding abscissas.
Higher homogeneous terms have larger total fresh-amplitude degree
and cannot cancel the detected first term.

Now suppose local ranks are one or \(w\).
For a rank-one factor \(R(T)\), choose a nonzero polynomial entry
\(h=R_{ab}\) and use
\[
 R(T)=(R(T)e_b)\bigl(e_a^{\mathsf T}R(T)/h(T)\bigr).
\]
Splitting at all rank-one factors gives nonzero rational scalar
segments with invertible middle runs.
Keep every denominator \(h\) outside its scalar numerator.
The new row endpoint \(e_a^{\mathsf T}R\) and column endpoint
\(Re_b\) are polynomial vectors with entry degree at most \(\delta\).
Complete each by constant standard basis rows or columns to a
generically invertible \(w\)-by-\(w\) polynomial matrix.
Its entry-degree bound is still \(\delta\).
Constant endpoint vectors admit constant completions.

Thus every scalar numerator is an invertible-local polynomial word
of the same local degree bound and is detected by the same fresh
\(\SV_{w-1}\). Each denominator remains nonzero, since its
constant fresh-amplitude value is \(h(C_i(t))\ne0\).
Neighboring segments may share a boundary variable; within each
segment every original variable occurs once.
The product of their nonzero rational substitutions is nonzero.
This proves the first assertion of Theorem~\ref{thm:hit}.

For pure \(\SV_w\), use its first pair as the initial shift
\[
 C_i=c_i+z_0\ell_i(\lambda(y_0)).
\]
The coefficient \(\ell_i(\lambda(y_0))\) is nonzero, since the
Lagrange polynomials are independent.
Over \(\F(y_0)\), each \(C_i\) is nonconstant in \(z_0\) and
therefore transcendental. The same argument applies and uses
the remaining \(w-1\) pairs as the fresh map.

Finally \(\deg f\le D=n\delta\).
For the explicit \(\gamma\), the individual seed degrees are at most
\(nD\) in \(t\), \((n-1)D\) in each direction, and \(D\)
in each amplitude. These are the rectangular grid bounds in
Theorem~\ref{thm:hit}. The finite evaluations use the original
polynomial, so numerical grid points need not avoid auxiliary
rational pivots separately.

\section{Small-characteristic boundary}

The derivative plateau proof does not justify an unrestricted
small-characteristic conclusion. For example, ordinary derivatives
can miss \(T^p\).
One may instead obtain latent coefficient concentration measured by
base-\(p\) digit weights using Hasse derivatives, but that alone
does not survive an arbitrary affine change as a low-support
coefficient assertion.

The present theorem uses ordinary total degree: once a nonzero
coefficient has degree at most \(w-1\), the entire first homogeneous
term has that ordinary degree bound. This is preserved by invertible
linear substitution. A low-digit-weight witness, possibly of much
higher ordinary degree than other terms, would not provide this
argument. No all-characteristic \(\gamma+\SV_{w-1}\) theorem for
general width is claimed here.
\end{document}
